\documentclass[varwidth=14cm,border=10pt]{standalone}
\usepackage{amsmath,xcolor,array}
\usepackage{fontspec,unicode-math}
\setmainfont{DejaVu Sans}[ItalicFont={DejaVu Sans},ItalicFeatures={FakeSlant=0.2},BoldItalicFont={DejaVu Sans Bold},BoldItalicFeatures={FakeSlant=0.2}]
\setmathfont{texgyredejavu-math.otf}
\setmathfont{DejaVu Sans}[range={up,bfup}]
\setmathfont{DejaVu Sans}[range={it,bfit},FakeSlant=0.2]
\setmathfont{texgyredejavu-math.otf}[range={"2202,"1D715,"03B5,"03F5,"2112,"2131}]% keep partial, epsilon, script caps
\pagecolor[HTML]{1E1E2E}\color[HTML]{F4F4F4}
\definecolor{accent}{HTML}{89B4FA}
\definecolor{warn}{HTML}{F38BA8}
\definecolor{good}{HTML}{A6E3A1}
\newcommand{\heading}[1]{\par\medskip{\large\bfseries\color{accent}#1}\par\smallskip}
\setlength{\parindent}{0pt}
\raggedright
\begin{document}

{\Large\bfseries L05 — Implicit sensitivity: FD, direct, adjoint}

\heading{1. The setting}
The design is $x\in\mathbb{R}^{n}$. The physics is a linear system whose matrix is built from
the design, linearly:
\[
  A(x)\,p = q,\qquad A(x)=\sum_{e=1}^{n} x_e A_e,\qquad \frac{\partial A}{\partial x_e}=A_e .
\]
The response is a linear functional of the state:
\[
  f(x) = c^{\!\top} p(x).
\]
There is {\itshape no formula} for $p$. It exists only as the output of a solve, so $f$ is an
{\bfseries implicit} function of $x$.

\heading{2. Where the implicitness lives}
In general $f$ can depend on $x$ twice over, and the chain rule keeps both:
\[
  \frac{df}{dx_e}
  = \underbrace{\frac{\partial f}{\partial x_e}\bigg|_{p\ \text{fixed}}}_{\text{explicit — free}}
  \; + \;
  \underbrace{\left(\frac{\partial f}{\partial p}\right)^{\!\top}\frac{\partial p}{\partial x_e}}_{\text{implicit — the hard one}} .
\]
For $f=c^{\!\top}p$ the explicit term vanishes ($f$ contains no $x$ of its own), so
$\partial f/\partial p = c$ and everything hangs on $\partial p/\partial x_e$.

\textcolor{good}{Do not over-generalise:} the volume budget $\sum_e L_e x_e \le V$ is the
opposite extreme — purely explicit, no state in it, gradient $=L_e$, free.

\heading{3. Method 1 — finite differences}
\[
  \frac{df}{dx_e}\approx\frac{f(x+h e_e)-f(x)}{h}
  \quad\text{(forward, } O(h)\text{)},\qquad
  \frac{f(x+h e_e)-f(x-h e_e)}{2h}
  \quad\text{(central, } O(h^2)\text{)}.
\]
\[
  \text{cost: } n+1 \text{ solves (forward)},\qquad 2n \text{ solves (central)} .
\]
Two errors fight: truncation $\sim h^{k}$ falls as $h\to0$, round-off
$\sim \varepsilon_{\text{mach}}/h$ grows. The optimal $h$ cannot be found without knowing the
answer. \textcolor{warn}{Verification only — never production.}

\heading{4. Method 2 — direct differentiation}
$A(x)p(x)=q$ holds for all $x$, so differentiate the whole equation ($q$ fixed):
\[
  \frac{\partial A}{\partial x_e}p + A\frac{\partial p}{\partial x_e} = 0
  \;\Longrightarrow\;
  \boxed{\;A\,\frac{\partial p}{\partial x_e} = -A_e\,p\;}
\]
Same matrix $A$, new right-hand side. $-A_e p$ is the {\bfseries pseudo-load}: the
out-of-balance you create by changing $x_e$ with the state frozen. Then
$df/dx_e = c^{\!\top}(\partial p/\partial x_e)$.
\[
  \text{cost: } 1 + n \text{ solves} \qquad\text{— one per \textcolor{accent}{design variable}.}
\]

\heading{5. Method 3 — the adjoint}
Attach a multiplier to the analysis, exactly as in L04:
\[
  \mathcal{L} = c^{\!\top}p + \lambda^{\!\top}\!\left(Ap-q\right).
\]
The bracket is identically zero at the solution, so $\mathcal{L}=f$ for {\itshape any}
$\lambda$. That freedom is the whole trick. Differentiating,
\[
  \frac{df}{dx_e}
  = \underbrace{\left(c^{\!\top}+\lambda^{\!\top}A\right)}_{\text{choose }\lambda\text{ to kill this}}
    \frac{\partial p}{\partial x_e}
  \;+\;\lambda^{\!\top}A_e\,p .
\]
Demand $c^{\!\top}+\lambda^{\!\top}A=0$:
\[
  \boxed{\;A^{\!\top}\lambda = -c\;}\qquad\Longrightarrow\qquad
  \boxed{\;\frac{df}{dx_e} = \lambda^{\!\top}A_e\,p\quad\text{for every }e\;}
\]
\[
  \text{cost: } 1 + n_r \text{ solves} \qquad\text{— one per \textcolor{accent}{response}.}
\]

\heading{6. The counting rule}
\begin{center}
\begin{tabular}{lll}
 & solves & indexed by\\[2pt]
direct  & $1+n$   & design variables\\
adjoint & $1+n_r$ & responses\\
\end{tabular}
\end{center}
Neither is better. Count $n$ against $n_r$, pick the smaller. One direct solve fills a whole
{\itshape column} of the Jacobian $\partial p_k/\partial x_e$; one adjoint solve fills a whole
{\itshape row}.

\heading{7. Self-adjointness}
If $A=A^{\!\top}$ and $c=q$ (the response is measured where the load is applied), then
$A\lambda=-q$ is the analysis already solved:
\[
  \lambda=-p\qquad\Longrightarrow\qquad
  \frac{df}{dx_e} = -\,p^{\!\top}A_e\,p \;\le\; 0
  \quad\text{if } A_e \succeq 0 .
\]
The adjoint solve disappears: the gradient is free. For the pipe network
$A_e=(e_i-e_j)(e_i-e_j)^{\!\top}$, so
\[
  \frac{df}{dx_e} = -\left(p_i-p_j\right)^{2}
  \quad\text{— minus the squared pressure drop across pipe } e .
\]

\end{document}
